\documentclass[10pt,a4paper]{amsart}
\usepackage{amsmath,amssymb,amsthm,graphicx}
% Retained minimal preamble for the Stacks Project chapter more-morphisms.tex (More on Morphisms).
% Verbatim excerpt of the chapter preamble: Xy-pic setup (native preamble lines 14-18) and the native
% theorem declarations with the mathematical macros (native preamble lines 154-217). Copied mechanically;
% nothing added, removed or redefined. The wrapper class and the amsmath/amssymb/amsthm/graphicx base
% packages are supplied by the assembler, so the class/externaldocument part of the native preamble is omitted.
% For commutative diagrams we use Xy-pic
\usepackage[all]{xy}

% We use 2cell for 2-commutative diagrams.
\xyoption{2cell}
% Theorem environments.
%
\theoremstyle{plain}
\newtheorem{theorem}[subsection]{Theorem}
\newtheorem{proposition}[subsection]{Proposition}
\newtheorem{lemma}[subsection]{Lemma}

\theoremstyle{definition}
\newtheorem{definition}[subsection]{Definition}
\newtheorem{example}[subsection]{Example}
\newtheorem{exercise}[subsection]{Exercise}
\newtheorem{situation}[subsection]{Situation}

\theoremstyle{remark}
\newtheorem{remark}[subsection]{Remark}
\newtheorem{remarks}[subsection]{Remarks}

\numberwithin{equation}{subsection}

% Macros
%
\def\lim{\mathop{\mathrm{lim}}\nolimits}
\def\colim{\mathop{\mathrm{colim}}\nolimits}
\def\Spec{\mathop{\mathrm{Spec}}}
\def\Hom{\mathop{\mathrm{Hom}}\nolimits}
\def\Ext{\mathop{\mathrm{Ext}}\nolimits}
\def\SheafHom{\mathop{\mathcal{H}\!\mathit{om}}\nolimits}
\def\SheafExt{\mathop{\mathcal{E}\!\mathit{xt}}\nolimits}
\def\Sch{\mathit{Sch}}
\def\Mor{\mathop{\mathrm{Mor}}\nolimits}
\def\Ob{\mathop{\mathrm{Ob}}\nolimits}
\def\Sh{\mathop{\mathit{Sh}}\nolimits}
\def\NL{\mathop{N\!L}\nolimits}
\def\CH{\mathop{\mathrm{CH}}\nolimits}
\def\proetale{{pro\text{-}\acute{e}tale}}
\def\etale{{\acute{e}tale}}
\def\QCoh{\mathit{QCoh}}
\def\Ker{\mathop{\mathrm{Ker}}}
\def\Im{\mathop{\mathrm{Im}}}
\def\Coker{\mathop{\mathrm{Coker}}}
\def\Coim{\mathop{\mathrm{Coim}}}

% Boxtimes
%
\DeclareMathSymbol{\boxtimes}{\mathbin}{AMSa}{"02}

%
% Macros for moduli stacks/spaces
%
\def\QCohstack{\mathcal{QC}\!\mathit{oh}}
\def\Cohstack{\mathcal{C}\!\mathit{oh}}
\def\Spacesstack{\mathcal{S}\!\mathit{paces}}
\def\Quotfunctor{\mathrm{Quot}}
\def\Hilbfunctor{\mathrm{Hilb}}
\def\Curvesstack{\mathcal{C}\!\mathit{urves}}
\def\Polarizedstack{\mathcal{P}\!\mathit{olarized}}
\def\Complexesstack{\mathcal{C}\!\mathit{omplexes}}
% \Pic is the operator that assigns to X its picard group, usage \Pic(X)
% \Picardstack_{X/B} denotes the Picard stack of X over B
% \Picardfunctor_{X/B} denotes the Picard functor of X over B
\def\Pic{\mathop{\mathrm{Pic}}\nolimits}
\def\Picardstack{\mathcal{P}\!\mathit{ic}}
\def\Picardfunctor{\mathrm{Pic}}
\def\Deformationcategory{\mathcal{D}\!\mathit{ef}}
\usepackage{cleveref}
\crefname{originallemma}{original lemma}{original lemmas}
\Crefname{originallemma}{Original lemma}{Original lemmas}
\makeatletter
\newlabel{cohomology-lemma-projection-formula}{{cohomology-lemma-projection-formula}{source}{}{sourcepointer-1}{}}
\newlabel{cohomology-lemma-projection-formula@cref}{{[originallemma][0][]cohomology-lemma-projection-formula}{[1][0][]source}{}{}{}}
\makeatother
\crefname{originallemma}{original lemma}{original lemmas}
\Crefname{originallemma}{Original lemma}{Original lemmas}
\makeatletter
\newlabel{morphisms-lemma-fppf-open}{{morphisms-lemma-fppf-open}{source}{}{sourcepointer-2}{}}
\newlabel{morphisms-lemma-fppf-open@cref}{{[originallemma][0][]morphisms-lemma-fppf-open}{[1][0][]source}{}{}{}}
\makeatother
\crefname{originallemma}{original lemma}{original lemmas}
\Crefname{originallemma}{Original lemma}{Original lemmas}
\makeatletter
\newlabel{perfect-lemma-flat-proper-perfect-direct-image-general}{{perfect-lemma-flat-proper-perfect-direct-image-general}{source}{}{sourcepointer-3}{}}
\newlabel{perfect-lemma-flat-proper-perfect-direct-image-general@cref}{{[originallemma][0][]perfect-lemma-flat-proper-perfect-direct-image-general}{[1][0][]source}{}{}{}}
\makeatother
\crefname{originallemma}{original lemma}{original lemmas}
\Crefname{originallemma}{Original lemma}{Original lemmas}
\makeatletter
\newlabel{perfect-lemma-locally-closed-where-H0-locally-free}{{perfect-lemma-locally-closed-where-H0-locally-free}{source}{}{sourcepointer-4}{}}
\newlabel{perfect-lemma-locally-closed-where-H0-locally-free@cref}{{[originallemma][0][]perfect-lemma-locally-closed-where-H0-locally-free}{[1][0][]source}{}{}{}}
\makeatother
\title{The diagonal of the Picard functor of a flat proper family is a locally closed immersion}
\author{The Stacks Project Authors}
\date{stacks-project:more-morphisms.tex:local-snapshot-20261011; GNU FDL 1.2 or later}
\begin{document}
\maketitle
\noindent\textit{Editorial note.} Statement span is lines 9198-9215 of sources/stacks-more-morphisms/pinned-original-source.tex; proof span is lines 9217-9304, the complete author proof (86 body lines).\par\medskip
\noindent\textit{Editorial note.} Navigation stubs: 4 cited-but-not-included labels are supplied as non-load-bearing pointers with their pinned tags where the snapshot tag index has them.\par\medskip
\noindent\textit{Editorial note.} extra\_packages is empty on purpose: the retained preamble already loads xy with options and a second declaration clashes in the assembler wrapper.\par\medskip
\noindent\textit{Editorial note.} Text is retained verbatim, including the author's own typos and placeholders.\par\medskip
\noindent\textit{Allowed standard prerequisites (editorial).}\begin{itemize}
\item Explicit prerequisite boundary. The Picard functor and rigidified line bundles, cohomology and base change for flat proper families, and supports of coherent modules; the locally-closedness argument is retained.
\end{itemize}\medskip
\ifdefined\hypertarget\hypertarget{sourcepointer-1}{}\fi
\noindent\textit{Source reference.} Original lemma cohomology-lemma-projection-formula: sources/stacks-more-morphisms/pinned-original-source.tex, source lines 9304--9304. Cited by the retained statement or proof but not part of the collected spans; supplied only as a navigation pointer. The argument does not depend on its content.\par\medskip
\ifdefined\hypertarget\hypertarget{sourcepointer-2}{}\fi
\noindent\textit{Source reference.} Original lemma morphisms-lemma-fppf-open: sources/stacks-more-morphisms/pinned-original-source.tex, source lines 9304--9304. Cited by the retained statement or proof but not part of the collected spans; supplied only as a navigation pointer. The argument does not depend on its content.\par\medskip
\ifdefined\hypertarget\hypertarget{sourcepointer-3}{}\fi
\noindent\textit{Source reference.} Original lemma perfect-lemma-flat-proper-perfect-direct-image-general: sources/stacks-more-morphisms/pinned-original-source.tex, source lines 9304--9304. Cited by the retained statement or proof but not part of the collected spans; supplied only as a navigation pointer. The argument does not depend on its content.\par\medskip
\ifdefined\hypertarget\hypertarget{sourcepointer-4}{}\fi
\noindent\textit{Source reference.} Original lemma perfect-lemma-locally-closed-where-H0-locally-free: sources/stacks-more-morphisms/pinned-original-source.tex, source lines 9304--9304. Cited by the retained statement or proof but not part of the collected spans; supplied only as a navigation pointer. The argument does not depend on its content.\par\medskip
\begin{lemma}
\label{lemma-diagonal-picard-flat-proper}
Let $f : X \to S$ be a flat, proper morphism of finite presentation.
Let $\mathcal{E}$ be a finite locally free $\mathcal{O}_X$-module.
For a morphism $g : T \to S$ consider the base change diagram
$$
\xymatrix{
X_T \ar[d]_p \ar[r]_q & X \ar[d]^f \\
T \ar[r]^g & S
}
$$
Assume $\mathcal{O}_T \to p_*\mathcal{O}_{X_T}$ is an
isomorphism for all $g : T \to S$. Then there exists an
immersion $j : Z \to S$ of finite presentation such that
a morphism $g : T \to S$ factors through $Z$ if and only if
there exists a finite locally free $\mathcal{O}_T$-module $\mathcal{N}$
with $p^*\mathcal{N} \cong q^*\mathcal{E}$.
\end{lemma}

\begin{proof}
Observe that the fibres $X_s$ of $f$ are connected by our assumption
that $H^0(X_s, \mathcal{O}_{X_s}) = \kappa(s)$. Thus the rank of
$\mathcal{E}$ is constant on the fibres. Since $f$ is open
(Morphisms, Lemma \ref{morphisms-lemma-fppf-open}) and closed
we conclude that there is a decomposition $S = \coprod S_r$
of $S$ into open and closed subschemes such that $\mathcal{E}$
has constant rank $r$ on the inverse image of $S_r$.
Thus we may assume $\mathcal{E}$ has constant rank $r$.
We will denote $\mathcal{E}^\vee = \SheafHom(\mathcal{E}, \mathcal{O}_X)$
the dual rank $r$ module.

\medskip\noindent
By cohomology and base change (more precisely by
Derived Categories of Schemes, Lemma
\ref{perfect-lemma-flat-proper-perfect-direct-image-general})
we see that $E = Rf_*\mathcal{E}$ is a perfect object of the
derived category of $S$ and that its formation commutes with
arbitrary change of base. Similarly for $E' = Rf_*\mathcal{E}^\vee$.
Since there is never any cohomology in degrees $< 0$, we see that
$E$ and $E'$ have (locally) tor-amplitude in $[0, b]$ for some $b$.
Observe that for any $g : T \to S$ we have
$p_*(q^*\mathcal{E}) = H^0(Lg^*E)$ and
$p_*(q^*\mathcal{E}^\vee) = H^0(Lg^*E')$.
Let $j : Z \to S$ and $j' : Z' \to S$ be immersions
of finite presentation constructed in Derived Categories of Schemes, Lemma
\ref{perfect-lemma-locally-closed-where-H0-locally-free}
for $E$ and $E'$ with $a = 0$ and $r = r$; these are roughly speaking
characterized by the property that $H^0(Lj^*E)$ and $H^0((j')^*E')$
are finite locally free modules compatible with pullback.

\medskip\noindent
Let $g : T \to S$ be a morphism. If there exists an $\mathcal{N}$
as in the lemma, then, using the projection formula
Cohomology, Lemma \ref{cohomology-lemma-projection-formula},
we see that the modules
$$
p_*(q^*\mathcal{E}) \cong
p_*(p^*\mathcal{N}) \cong
\mathcal{N} \otimes_{\mathcal{O}_T} p_*\mathcal{O}_{X_T} \cong
\mathcal{N}\quad\text{and similarly }\quad
p_*(q^*\mathcal{E}^\vee) \cong \mathcal{N}^\vee
$$
are finite locally free modules of rank $r$
and remain so after any further base change $T' \to T$.
Hence in this case $T \to S$ factors through $j$ and through $j'$.
Thus we may replace $S$ by $Z \times_S Z'$ and assume that
$f_*\mathcal{E}$ and $f_*\mathcal{E}^\vee$ are finite locally free
$\mathcal{O}_S$-modules of rank $r$
whose formation commutes with arbitrary change of base
(small detail omitted).

\medskip\noindent
In this situation if $g : T \to S$ be a morphism and there exists an
$\mathcal{N}$ as in the lemma, then the map (cup product in degree $0$)
$$
p_*(q^*\mathcal{E})
\otimes_{\mathcal{O}_T}
p_*(q^*\mathcal{E}^\vee)
\longrightarrow \mathcal{O}_T
$$
is a perfect pairing. Conversely, if this cup product map is a
perfect pairing, then we see that locally on $T$ we may choose a
basis of sections
$\sigma_1, \ldots, \sigma_r$ in $p_*(q^*\mathcal{E})$
and $\tau_1, \ldots, \tau_r$ in $p_*(q^*\mathcal{E}^\vee)$
whose products satisfy $\sigma_i \tau_j = \delta_{ij}$.
Thinking of $\sigma_i$ as a section of $q^*\mathcal{E}$ on $X_T$
and $\tau_j$ as a section of $q^*\mathcal{E}^\vee$ on $X_T$,
we conclude that
$$
\sigma_1, \ldots, \sigma_r :
\mathcal{O}_{X_T}^{\oplus r}
\longrightarrow
q^*\mathcal{E}
$$
is an isomorphism with inverse given by
$$
\tau_1, \ldots, \tau_r :
q^*\mathcal{E}
\longrightarrow
\mathcal{O}_{X_T}^{\oplus r}
$$
In other words, we see that $p^*p_*q^*\mathcal{E} \cong q^*\mathcal{E}$.
But the condition that the cup product is nondegenerate picks
out a retrocompact open subscheme (namely, the locus where a suitable
determinant is nonzero) and the proof is complete.
\end{proof}

\end{document}
